% !PRINTED 71
\weberpage{71}
%<<CONT
obtains. But this is the group
\weq{7}{\mathfrak{A}_\mu \equiv 1\,(\mathrm{mod}\ \mu),}

that is, the group consisting of all those numbers of $\mathfrak{N}$ which satisfy
the congruence $a \equiv 1\,(\mathrm{mod}\ \mu)$. For first of all it is clear
that $\mathfrak{A}_{m_1}$ and $\mathfrak{A}_{m_2}$ are divisors of $\mathfrak{A}_\mu$. Hence the
least common multiple of $\mathfrak{A}_{m_1}$ and $\mathfrak{A}_{m_2}$ is also contained in $\mathfrak{A}_\mu$.
But if, on the other hand, $\alpha = 1 + \xi\mu$ is any number in $\mathfrak{A}_\mu$, then
one can determine the numbers $a_1 = 1 + x_1 m_1$ in $\mathfrak{A}_{m_1}$ and $a_2 = 1 + x_2 m_2$
in $\mathfrak{A}_{m_2}$ in such a way that $\alpha \equiv a_1 a_2\,(\mathrm{mod}\ m)$. One has
only to set $\xi\mu = x_1 m_1 + x_2 m_2$, which by Bd. I, §. 118 is always
possible. Hence $\mathfrak{A}_\mu$ is also contained in the least common multiple of $\mathfrak{A}_{m_1}$ and $\mathfrak{A}_{m_2}$, which is accordingly
$\mathfrak{A}_\mu$ itself. From this, then, the theorem:

\gesperrt{4. If $m_1$ and $m_2$ are any two natural numbers
and $\mu$ is the greatest common divisor of
$m_1$ and $m_2$, then the field $\Omega_\mu$ is the greatest
common divisor of $\Omega_{m_1}$ and $\Omega_{m_2}$, and
$\Omega_{\frac{m_1 m_2}{\mu}}$ their least common multiple.}

From this it now follows: if two full cyclotomic fields $\Omega_m$
and $\Omega_{m'}$ have a common divisor $\Omega$, and $m'$ is smaller
than $m$, then there must be a proper divisor $d$ of $m$ such that
$\Omega$ is also a divisor of $\Omega_d$.

We shall call a divisor of $\Omega_m$ \gesperrt{primary} if it is
not at the same time contained in a full cyclotomic field $\Omega_{m'}$ of
lower $m'$. It then follows that one obtains all
the non-primary divisors of $\Omega_m$ by letting the
index $d$ in $\Omega_d$ run through all the proper divisors of $m$ and
seeking out the divisors of $\Omega_d$.

If now we denote by $q_1, q_2, \ldots q_\tau$ all the distinct
primes dividing $m$ (2 included) and
set
\weq{8}{m = q_1 m_1 = q_2 m_2 = \cdots = q_\tau m_\tau,}

then every $d$ is a divisor of one of $m_1, m_2, \ldots m_\tau$, and we obtain the non-primary divisors of $\Omega_m$ if, for the fields
\weq{9}{\Omega_{m_1},\ \Omega_{m_2},\ \ldots\ \Omega_{m_\tau}}
we seek out all their divisors.

These fields, however, belong to the group
\weq{10}{Q_1 \equiv 1\,(\mathrm{mod}\ m_1),\ \ Q_2 \equiv 1\,(\mathrm{mod}\ m_2),\ \ldots\ Q_\tau \equiv 1\,(\mathrm{mod}\ m_\tau),}
%CONT>>
